\documentclass[11pt]{article}

\usepackage[a4paper,margin=25mm]{geometry}
\usepackage[T1]{fontenc}
\usepackage[utf8]{inputenc}
\usepackage{lmodern}
\usepackage{microtype}
\usepackage{amsmath,amssymb,mathtools}
\usepackage{amsthm}
\usepackage{array}
\usepackage{booktabs}
\usepackage{enumitem}
\usepackage{tabularx}
\usepackage{xcolor}
\usepackage{hyperref}

\definecolor{artifactblue}{HTML}{1F4E79}
\definecolor{artifactgreen}{HTML}{1E6B45}
\definecolor{artifactred}{HTML}{8B2F2F}

\hypersetup{
  colorlinks=true,
  linkcolor=artifactblue,
  urlcolor=artifactblue,
  pdftitle={A Mathematical Reading of TargetKeygenEtaSampler.ec},
  pdfsubject={HAETAE centered-trit consumer correctness}
}

\setlist[itemize]{leftmargin=1.5em,itemsep=0.25em,topsep=0.35em}
\setlist[enumerate]{leftmargin=1.7em,itemsep=0.25em,topsep=0.35em}
\renewcommand{\arraystretch}{1.18}
\setlength{\emergencystretch}{3em}

\theoremstyle{definition}
\newtheorem{definition}{Definition}
\theoremstyle{plain}
\newtheorem{theorem}{Theorem}
\newtheorem{proposition}{Proposition}
\theoremstyle{remark}
\newtheorem*{remark}{Remark}

\newcommand{\code}[1]{\nolinkurl{#1}}
\newcommand{\Z}{\mathbb Z}
\newcommand{\uint}[1]{\langle #1\rangle_{\mathrm u}}
\newcommand{\sint}[1]{\langle #1\rangle_{\mathrm s}}
\newcommand{\Frame}{\mathsf{Frame}_{256}}
\newcommand{\WordFrame}{\mathsf{WordFrame}}
\newcommand{\Centered}{\mathsf{Centered}}

\newenvironment{resultbox}
  {\begin{center}\begin{minipage}{0.94\linewidth}\color{artifactgreen}\hrule\vspace{0.6em}\color{black}}
  {\vspace{0.6em}\color{artifactgreen}\hrule\end{minipage}\end{center}}

\newenvironment{boundarybox}
  {\begin{center}\begin{minipage}{0.94\linewidth}\color{artifactred}\hrule\vspace{0.6em}\color{black}}
  {\vspace{0.6em}\color{artifactred}\hrule\end{minipage}\end{center}}

\title{Centered-Trit Extraction for HAETAE\\
\large A mathematical reading of \code{TargetKeygenEtaSampler.ec}}
\author{HAETAE reference EasyCrypt artifact}
\date{Artifact state checked 2026-07-23}

\begin{document}
\maketitle

\begin{abstract}
The eta consumer reads bytes, rejects the thirteen values from $243$ through
$255$, and turns each accepted byte into as many as five coefficients in
$\{-1,0,1\}$.  The EasyCrypt development checks the machine arithmetic,
exact capped decoder, counter, centered prefix, and frame.  It also follows
successive exact 136-byte SHAKE256 blocks through the generated whole eta leaf
and lifts the finite decoded prefix through the actual eta vector caller.  The
finite consumer terminates unconditionally; the actual whole leaf terminates
with probability one under an explicit deterministic finite-prefix progress
certificate.  These facts do not prove that every input has such a certificate,
a probability distribution, or the security of HAETAE.
\end{abstract}

\begin{resultbox}
\textbf{The result in one sentence.}
Whenever the extracted consumer returns from a state satisfying the stated
capacity and prefix assumptions, it has written at most 256 coefficients,
every coefficient in its established prefix represents $-1$, $0$, or $1$,
and every 32-bit word outside the reserved polynomial block is unchanged.  The
finite consumer is unconditionally lossless, and the actual whole eta leaf is
lossless under the explicit \code{eta_progress_prefix} certificate and its
exact seed, slice, base, and capacity premises.
\end{resultbox}

\section{The Mathematical Question}

The procedure under study is
\begin{center}
\code{__kp_poly_uniform_eta_consume_2048}.
\end{center}
It takes an output array $A$, a base $b$, a coefficient counter $k$, a byte
buffer $B$, and a buffer length $L$.  It returns an updated pair $(A',k')$.
All five inputs are values in the extracted EasyCrypt model; there are no
physical pointers in the theorem statement.

The procedure is deterministic once the buffer is fixed.  Its local question
is therefore not ``is this a random sampler?'' but rather:
\begin{quote}
Which values can it write, how far can its counter advance, and which parts of
the output array can it modify?
\end{quote}

The authored file \code{KeygenEtaSamplerSpec.ec} supplies mathematical
predicates and arithmetic lemmas.  The file
\code{TargetKeygenEtaSampler.ec} applies them directly to the generated
procedure through ordinary Hoare judgments.  It does not compare the target
with a separate high-level sampling procedure.

\section{Words, Arrays, and Polynomial Blocks}

A 32-bit word $w$ has an unsigned value
$\uint{w}\in\{0,\ldots,2^{32}-1\}$ and a signed two's-complement value
$\sint{w}\in\{-2^{31},\ldots,2^{31}-1\}$.  A 64-bit word has the analogous
unsigned interpretation; the intended interpretation will always be clear
from context.

The output type \code{BArray8192} contains $8192$ bytes.  The proof accesses
it in aligned 32-bit units, so its word capacity is
\[
  C=8192/4=2048.
\]
Write $A[i]$ for the 32-bit word returned by
\code{BArray8192.get32 A i}.  A polynomial occupies
\[
  N=256
\]
consecutive words.  Thus a base $b$ is a word index, not a byte offset.  The
input buffer \code{BArray1024} contains $1024$ bytes and is read with
\code{get8}.

\begin{definition}[Frame and centeredness predicates]
For arrays $A_0,A$, an integer base $b$, and integers $s,t$, define
\begin{align*}
 \Frame(A_0,A;b)
   &\iff \forall i\in[0,C),\quad
      i\notin[b,b+N)\Longrightarrow A[i]=A_0[i],\\
 \WordFrame_{s:t}(A_0,A;b)
   &\iff \forall i\in[0,C),\quad
      i\notin[b+s,b+t)\Longrightarrow A[i]=A_0[i],\\
 \Centered(A;b,s,t)
   &\iff \forall i\in[s,t),\quad
      -1\le \sint{A[b+i]}\le1.
\end{align*}
\end{definition}

The first predicate permits changes anywhere in the complete reserved block.
The second is a tighter, interval-specific frame predicate.  The specification
proves useful one-write lemmas about both predicates, but the final consumer
theorem uses only the coarse predicate $\Frame$.  In particular, the authored
theorem does not state a tight frame for exactly the newly filled interval.

The centered predicate is also deliberately modest: it records only the
range $\{-1,0,1\}$.  It does not remember which input byte produced a
coefficient or assert any probability law.

\section{The Centered Residue Map}

For a 32-bit word $x$, define
\[
  r(x)=\uint{x}\bmod 3
\]
and
\[
  c(x)=
  \begin{cases}
    -1,&r(x)=2,\\
    r(x),&r(x)=0\text{ or }1.
  \end{cases}
\]
The word written by the specification is
\[
  \widehat c(x)=\operatorname{W32.of\_int}(c(x)).
\]
The representation of $-1$ is therefore the word $2^{32}-1$, whose signed
interpretation is $-1$.

\begin{proposition}[Centered residue facts]
For every 32-bit word $x$,
\[
  0\le r(x)<3,
  \qquad -1\le c(x)\le1,
  \qquad \sint{\widehat c(x)}=c(x).
\]
Consequently $-1\le\sint{\widehat c(x)}\le1$.
\end{proposition}

The generated helper names contain \code{mod3}, but they return this
\emph{centered} residue rather than the ordinary residue in $\{0,1,2\}$.

\section{Five Base-3 Digits from One Byte}

The implementation replaces division by $3$ with a multiplication and a
right shift.  The key integer identity is
\begin{equation}
  \left\lfloor\frac{171x}{512}\right\rfloor
  =\left\lfloor\frac{x}{3}\right\rfloor,
  \qquad 0\le x\le242.
  \label{eq:div3}
\end{equation}
Since $512=2^9$, the left side is implemented by multiplying by $171$ and
shifting right by nine bits.  The proof also establishes the bound cascade
\[
  242\longmapsto80\longmapsto26\longmapsto8\longmapsto2.
\]
These small bounds justify every later machine-word calculation and exclude
overflow in the multiplication used here.

For an accepted byte $x<243=3^5$, put
\[
  q_j=\left\lfloor\frac{x}{3^j}\right\rfloor,
  \qquad
  d_j=q_j\bmod3,
  \qquad 0\le j<5.
\]
Then
\[
  x=\sum_{j=0}^{4}d_j3^j,
\]
and the five intended outputs are
\[
  c(q_0),\ c(q_1),\ c(q_2),\ c(q_3),\ c(q_4).
\]
The last quotient satisfies $0\le q_4\le2$.  The generated code centers it
with the formula
\[
  q_4-3\left\lfloor\frac{q_4}{2}\right\rfloor,
\]
which gives $0,1,-1$ for $q_4=0,1,2$, respectively.

Let $M_{242}$, $M_{26}$, and $M_8$ denote, respectively, the generated
procedures \code{__poly_sample_mod3},
\code{__poly_sample_mod3_leq26}, and
\code{__poly_sample_mod3_leq8}.  The refinement file proves the exact
contracts
\begin{align*}
 \uint{x}\le242
   &\Longrightarrow M_{242}(x)=\widehat c(x),\\
 \uint{x}\le26
   &\Longrightarrow M_{26}(x)=\widehat c(x),\\
 \uint{x}\le8
   &\Longrightarrow M_8(x)=\widehat c(x).
\end{align*}

\section{A Mathematical Transcription of the Consumer}
\label{sec:transcription}

For an ordinary in-bounds buffer length $0\le L\le1024$, the generated
control flow has the following transparent form.  This is a transcription of
the code, not an exact-output theorem stated in the current EasyCrypt files.

\begin{quote}
\begin{verbatim}
p := 0
while k < 256 and p < L:
    x := unsigned value of B[p]
    p := p + 1

    if x < 243:
        q := x
        for j = 0,1,2,3,4:
            if k = 256: break
            A[b+k] := encoded centered residue c(q)
            k := k + 1
            if j < 4: q := floor(q/3)

return (A,k)
\end{verbatim}
\end{quote}

Thus bytes $243,\ldots,255$ are rejected: consuming such a byte changes
neither $A$ nor $k$.  An accepted byte contributes up to five coefficients,
but every write after the first is guarded by $k<256$.  The first write is
safe because the outer loop has already established $k<256$.  The loop stops
when the polynomial is full or the declared buffer length is exhausted.

The actual target uses 64-bit words for $b,k,p,L$.  Its capacity hypothesis
implies that, while $k<256$,
\[
  \uint{b+k}=\uint{b}+\uint{k},
\]
so the machine address is the ordinary word index used above.

\section{The Three Machine-Checked Consumer Theorems}

An EasyCrypt Hoare judgment in this section has a partial-correctness reading:
if the procedure returns from a state satisfying the precondition, its result
satisfies the postcondition.  Separately, \code{eta_consume2048_ll} proves
unconditional termination of this finite consumer using a bounded variant;
that fact does not say that one block fills all 256 positions.

\begin{theorem}[Counter bound]
Let $k_0$ be the initial 64-bit counter.  If
\[
  k=k_0
  \qquad\text{and}\qquad
  \uint{k_0}\le256,
\]
then every returned counter $k'$ satisfies
\[
  \uint{k_0}\le\uint{k'}\le256.
\]
\end{theorem}

This theorem neither says that the counter reaches $256$ nor gives a formula
for the number of accepted or rejected bytes.

\begin{theorem}[Centered prefix]
Let $A_0$ be the initial array, $k_0$ the initial counter, and $b_0$ the
integer represented by the 64-bit base.  Suppose
\[
\begin{gathered}
 A=A_0,\qquad k=k_0,\qquad \uint{b}=b_0,\\
 b_0+256\le2048,\qquad \uint{k_0}\le256,\\
 \Centered(A_0;b_0,0,\uint{k_0}).
\end{gathered}
\]
If the consumer returns $(A',k')$, then
\[
  \Centered(A';b_0,0,\uint{k'}).
\]
\end{theorem}

The initial-prefix assumption allows the consumer to resume filling a
partially constructed polynomial.  The conclusion states only centeredness,
not equality of the old prefix or the origin of each new value.

\begin{theorem}[Reserved-block frame]
Let $A_0$ be the initial array and $b_0$ the integer represented by the
64-bit base.  Suppose
\[
  A=A_0,\qquad \uint{b}=b_0,\qquad
  b_0+256\le2048,\qquad 0\le\uint{k}\le256.
\]
If the consumer returns $(A',k')$, then
\[
  \Frame(A_0,A';b_0).
\]
Equivalently, every 32-bit word outside
$[b_0,b_0+256)$ has its original value.
\end{theorem}

These are three separately stated Hoare theorems about the same deterministic
procedure.  The source does not package their conclusions into one combined
theorem.

\section{Why the Proof Works}

The generated program contains one outer consumption loop.  The refinement
checks that loop three times, choosing the invariant appropriate to each
postcondition.

\subsection{Counter induction}

With initial counter $k_0$, the invariant is
\[
  \uint{k_0}\le\uint{k}\le256.
\]
Rejected bytes leave $k$ unchanged.  Each successful write increments $k$ by
one, and a write occurs only while $k<256$.  Hence the invariant is preserved.
This proof inlines the three centered-residue helpers because their returned
values are irrelevant to the counter argument.

\subsection{Centered-prefix induction}

The invariant is
\[
\begin{aligned}
 &\Centered(A;b_0,0,\uint{k}),\\
 &\uint{b}=b_0,\qquad b_0+256\le2048,\qquad
 0\le\uint{k}\le256.
\end{aligned}
\]
On a successful step, the helper-correctness theorem supplies a word whose
signed value is between $-1$ and $1$.  The program writes it at the next index
$b_0+\uint{k}$ and increments $k$.  The single-write centeredness lemma then
extends the interval from $[0,\uint{k})$ to $[0,\uint{k}+1)$.  This argument is
repeated for each of the five possible digits.

\subsection{Frame induction}

The invariant is
\[
\begin{aligned}
 &\Frame(A_0,A;b_0),\\
 &\uint{b}=b_0,\qquad b_0+256\le2048,\qquad
 0\le\uint{k}\le256.
\end{aligned}
\]
Here the value returned by a residue helper may be treated as arbitrary.  The
only relevant fact is that the next write index lies inside
$[b_0,b_0+256)$.  A single-write frame lemma therefore preserves the
invariant.  This separation is useful: value correctness proves centeredness,
whereas address correctness proves confinement.

The stronger block theorem records the exact emitted list as
$\operatorname{take}_{256}$ of the incoming prefix followed by the decoded
block.  The whole-leaf stream invariant composes those lists across successive
exact SHAKE256 blocks.

\section{Certificate-Conditioned Whole-Leaf Termination}

For fixed seed array $S$, seed offset $o$, nonce $\nu$, and block limit $L$,
let $C(r)$ be the size of the capped eta decoding of the first $r$ consecutive
136-byte SHAKE256 blocks from the exact padded state for
$S[o:o+64]\Vert\operatorname{LE}_2(\nu)$.  The predicate
\code{eta_progress_prefix S o nonce L}, written
$P_\eta(S,o,\nu,L)$ below, is a deterministic certificate with two parts:
\begin{enumerate}
\item $L\ge1$ and $C(L)=256$; and
\item for every $1\le r<L$, if $C(r)<256$, then $C(r)<C(r+1)$.
\end{enumerate}
It is a property of one concrete finite SHAKE256 byte sequence, not a
randomness assumption.  In particular, the strict-progress clause rules out
an all-rejected next block at every certified incomplete prefix.

\begin{theorem}[Actual eta-leaf probability-one termination]
Fix an initial seed array $S$, seed offset $o$, nonce $\nu$, word base $b$,
and limit $L$.  Suppose the generated leaf receives exactly those seed,
offset, nonce, and base values, and
\[
  b+256\le2048,
  \qquad \uint{o}+64\le128,
  \qquad P_\eta(S,o,\nu,L).
\]
Then
\code{_kp_poly_uniform_eta_at_seedbuf_2048} terminates with probability one.
\end{theorem}

This is the pHoare theorem \code{eta2048_leaf_progress_ll}.  Because the
generated leaf is deterministic, probability one here is a losslessness
statement under the displayed precondition.  The development does not prove
\code{eta_progress_prefix} for every seed, offset, and nonce; it also does not
lift per-leaf certificates through a parameterized eta-vector caller in this
file.  A separate parent refinement lifts five certificates through the two
fixed eta callers of the first attempt, but not through the actual mode-2
retry loop.

\section{The Distributional Bridge Is Conditional}

\begin{boundarybox}
\textbf{Conditional mathematical observation, not a theorem in these files.}
If $X$ is uniform on the 256 byte values, then
\[
  \Pr[X<243]=\frac{243}{256}.
\]
Conditional on $X<243$, the value $X$ is uniform on
$\{0,\ldots,242\}$.  Base-3 expansion is a bijection
\[
  \{0,\ldots,242\}\longleftrightarrow\{0,1,2\}^{5},
\]
so its five digits are jointly uniform.  Applying the bijection
$0\mapsto0$, $1\mapsto1$, $2\mapsto-1$ makes the five centered digits jointly
uniform on $\{-1,0,1\}^{5}$.
\end{boundarybox}

This observation explains the threshold $243=3^5$ and its relevance to an
eta-$1$ sampler.  To turn it into an implementation theorem one would still
need the appropriate distributional assumptions or game transition and a
proof that suitable progress certificates exist with the required
probability.  The deterministic buffer-to-SHAKE256 connection, refill
sequence, rejection decisions, 256-value cap, and certificate-conditioned
leaf termination are already machine checked.

\section{Security Relevance and Exact Boundary}

The proved facts are useful local obligations for a future HAETAE security
argument:
\begin{itemize}
\item centeredness establishes the support needed for coefficient and norm
      bounds;
\item the counter bound prevents more than 256 coefficients from being
      established in one polynomial block; and
\item the frame theorem isolates the block that the consumer may modify,
      which is necessary when composing several polynomial operations.
\end{itemize}

They stop before probabilistic or end-to-end security reasoning.

\begin{center}
\small
\begin{tabularx}{0.98\textwidth}{>{\raggedright\arraybackslash}p{0.57\textwidth}X}
\toprule
Claim & Status here \\
\midrule
Machine arithmetic returns centered residues on its stated domains & Proved \\
Returned counter lies between its initial value and $256$ & Proved \\
Established prefix contains only $-1,0,1$ & Proved \\
Words outside the reserved 256-word block are unchanged & Proved \\
Acceptance removes modulo bias for an independent uniform byte & Conditional observation \\
Input bytes are uniform or independent & Not proved \\
Finite leaf blocks equal the checked SHAKE256 stream & Proved on return \\
The finite consumer terminates & Proved unconditionally \\
One finite consumer call fills all 256 positions & Not guaranteed \\
The returned array equals an exact capped flattened digit sequence & Proved on return \\
The tight frame covers exactly the newly written interval & Not stated for the consumer \\
\code{buflen} is at most the 1024-byte buffer capacity & Not assumed here \\
The whole eta leaf has an exact finite-stream/range/frame contract & Proved on return \\
The whole eta leaf terminates under \code{eta_progress_prefix} and the exact bounds & Proved with probability one \\
Every seed and nonce has a suitable progress certificate & Not proved \\
The two fixed first-attempt eta callers terminate under five bundled
certificates & Proved in the separate parent refinement \\
The actual mode-2 retry loop terminates & Not proved \\
The full key generator is semantically composed & Not proved \\
Constant-time, speculative-execution, or leakage security & Not proved here \\
HAETAE security in a cryptographic game & Not proved here \\
\bottomrule
\end{tabularx}
\end{center}

In particular, the generated target contains erased declassification markers
and protection operations, but the judgments in this file do not use a
leakage semantics.  Also, the frame theorem compares 32-bit words; it is not
presented as a theorem about physical pointer separation or source-level
memory safety.

\section{Translation Back to the EasyCrypt Files}

The mathematical definitions and the extracted implementation are separated
as follows.

\noindent\begin{tabularx}{\textwidth}{>{\raggedright\arraybackslash}p{0.24\textwidth}X}
\toprule
Artifact & Mathematical role \\
\midrule
Authored specification & Centered-residue map, frame and interval predicates,
division identity, no-wrap fact, and one-write preservation lemmas. \\
Authored refinement & Helper arithmetic contracts and the three Hoare theorems
for the extracted consumer. \\
Generated target & Bit-level residue helpers and the concrete byte-consumption
loop transcribed in Section~\ref{sec:transcription}. \\
\bottomrule
\end{tabularx}

The corresponding paths are:
\begin{center}
\footnotesize
\code{easycrypt/spec/KeygenEtaSamplerSpec.ec}\\
\code{easycrypt/refinement/TargetKeygenEtaSampler.ec}\\
\code{easycrypt/extract/keygen-sampler-callers/KeygenSamplerCallersTarget.ec}
\end{center}

\begin{center}
\footnotesize
\begin{tabularx}{0.98\textwidth}{>{\raggedright\arraybackslash}p{0.39\textwidth}>{\raggedright\arraybackslash}X}
\toprule
Mathematical statement & EasyCrypt lemma \\
\midrule
$0\le r(x)<3$ & \code{base3_residue_range} \\
$-1\le c(x)\le1$ and correct signed encoding &
\code{centered_trit_value_range}, \code{centered_trit_to_sint} \\
Equation~\eqref{eq:div3} & \code{div3_via_171} \\
64-bit base-plus-counter does not wrap & \code{base_counter_no_wrap8192} \\
One write preserves the coarse frame & \code{poly_frame8192_set32_word} \\
One centered write extends the filled interval &
\code{centered_interval8192_set32_word} \\
Multiply-and-shift computes division by three & \code{div3_step_word} \\
The three centered-residue helper contracts &
\code{mod3_full_correct}, \code{mod3_leq26_correct},
\code{mod3_leq8_correct} \\
Counter theorem & \code{eta_consume2048_counter} \\
Centered-prefix theorem & \code{eta_consume2048_centered} \\
Reserved-block frame theorem & \code{eta_consume2048_frame} \\
Exact 136-byte consumer contract & \code{eta_consume2048_block136} \\
Unconditional finite-consumer termination & \code{eta_consume2048_ll} \\
Exact whole-leaf finite stream & \code{eta2048_leaf_stream} \\
Whole-leaf centered range and frame & \code{eta2048_leaf_correct} \\
Certificate-conditioned whole-leaf termination & \code{eta2048_leaf_progress_ll} \\
\bottomrule
\end{tabularx}
\end{center}

\section{Replaying the Machine-Checked Evidence}

From the \code{haetae-ref-easycrypt/} directory, run
\begin{center}
\code{./scripts/verify-keygen-sampler-callers-proof.sh}
\end{center}
This proof gate includes both \code{KeygenEtaSamplerSpec.ec} and
\code{TargetKeygenEtaSampler.ec}, together with their generated target and
the other sampler-caller proof units.  It checks source and extraction drift,
compiles the authored theories, and scans them for proof holes and authored
axiom declarations.

Build this guide independently with
\begin{quote}
\begin{verbatim}
cd haetae-ref-easycrypt/latex
pdflatex -interaction=nonstopmode -halt-on-error \
  -output-directory=build TargetKeygenEtaSampler.tex
pdflatex -interaction=nonstopmode -halt-on-error \
  -output-directory=build TargetKeygenEtaSampler.tex
\end{verbatim}
\end{quote}
The PDF is written to \code{build/TargetKeygenEtaSampler.pdf}.

\section*{Conclusion}

An accepted byte has five base-3 digits; verified word-level arithmetic maps
each digit into $\{-1,0,1\}$.  The proof establishes the exact capped decoder,
counter control, centeredness, confinement to one block, and its composition
with successive exact SHAKE256 blocks through the generated whole leaf.  The
bounded consumer terminates unconditionally, and the actual leaf terminates
with probability one under its explicit deterministic progress certificate.
Universal certificate existence, caller totality, and the probabilistic or
cryptographic sampler proof remain open in general; the two fixed
first-attempt eta callers and proof-only prefix now have
certificate-conditioned totality.

\end{document}
